\documentclass[11pt]{article}
\usepackage[a4paper,margin=26mm]{geometry}
\usepackage{amsmath,amssymb,amsthm,booktabs,hyperref}
\newtheorem{theorem}{Theorem}\newtheorem{lemma}[theorem]{Lemma}
\newcommand{\F}{\mathbb F}\newcommand{\gen}[1]{\langle#1\rangle}
\newcommand{\PSL}{\operatorname{PSL}}\newcommand{\PGL}{\operatorname{PGL}}
\title{Small regular subgroups and their translation kernels}
\author{A recovery of Sections 5.1--5.2 of Leemans--Vandenschrick}\date{}
\begin{document}\maketitle
Let $F=\F_q$ have characteristic $p$. We first recover the rank-two
order bound in Section~5.1, with its unipotent correction. For odd $p$
we then recover the reducible and conic rank-three alternatives of
Theorems~5.1--5.2. The actual translation modules also give corrected
types for the generic affine row of Table~3.
We retain the standard string end-intersection criterion, and Dickson's
classification of subgroups of $\PGL_2(F)$ in its natural subfield
form. The latter is an attributed classical input, not a new
classification claimed here.

\section*{Two involutions}
\begin{theorem}\label{two}
If distinct involutions generate a rank-two subgroup of $\PSL_3(q)$,
its order divides $2(q-1)$, $2(q+1)$ or $4p$. In characteristic two,
the last possibility has rotation order at most four. In odd
characteristic a regular unipotent rotation has order $p$.
\end{theorem}
\begin{proof}
The scalar centre has odd order dividing three. Each projective
involution has a unique involutory lift to determinant one, and the
product $X$ of the two lifts is conjugate to $X^{-1}$. Its eigenvalues
are therefore $1,\lambda,\lambda^{-1}$. If $\lambda\ne\pm1$, the
matrix is semisimple: either $\lambda\in F$ and its order divides
$q-1$, or $\lambda^q=\lambda^{-1}$ and its order divides $q+1$.
If all eigenvalues are one, the largest Jordan block has size at most
three. Its order divides $p$ for odd $p$, and four for $p=2$.
In odd characteristic the remaining case has eigenvalues $1,-1,-1$;
its order divides $2p$. Projectivization does not shorten any of these
orders: if $X^m=\zeta I$, its eigenvalue one forces $\zeta=1$.
The group generated by the two involutions is dihedral of twice the
rotation order, giving the assertion.

These are the expected ambient dihedral families, not just numerical
bounds. The usual split and nonsplit tori in $\PGL_2(q)$ have
normalizers $D_{2(q-1)}$ and $D_{2(q+1)}$; their three-dimensional
orthogonal representation embeds in $\PSL_3(q)$ when $q$ is odd.
When $q$ is even the block embedding of $\operatorname{SL}_2(q)$
gives the same normalizers. In odd characteristic,
$\operatorname{diag}(-1,-1,1)(I+E_{12})$ has order $2p$ and is inverted
by $\operatorname{diag}(1,-1,-1)$, giving $D_{4p}$.
In characteristic two, $I+E_{12}$ and $I+E_{23}$ give $D_8$.
Their dihedral subgroups realize the corresponding smaller orders.
\end{proof}
The source's displayed regular unipotent on physical page~11 is
incorrectly assigned the bound $2(2,p)$. Its order is $p$ for odd $p$
and four for $p=2$. The final dihedral order bound remains valid.
This is an isomorphism and order classification; no assertion that all
embeddings of a given dihedral group are conjugate is needed.

\section*{A common block form in odd characteristic}
Now suppose $p$ is odd. Put
\[
 d_1=\operatorname{diag}(1,-1,-1),\quad
 d_2=\operatorname{diag}(-1,1,-1),\quad d_3=d_1d_2.
\]
The commuting distinct end generators of a rank-three marking generate
the diagonal Klein group $D=\{1,d_1,d_2,d_3\}$ after a coordinate
change. The middle generator is a homology, with a centre and a
complementary axis. Suppose first that the whole group fixes a point
or line. The common eigenspaces of $D$ are exactly the coordinate
points and coordinate planes. After duality if necessary, the middle
axis therefore contains a coordinate point. Move this point to $e_3$.
The middle lift has the form
\[
 t=\begin{pmatrix}A&0\\v&-1\end{pmatrix},\qquad
 A=\begin{pmatrix}a&b\\c&-a\end{pmatrix},\quad
 a^2+bc=1,\quad vA=v.                                 \tag{1}
\]
The group is $G=\gen{D,t}$, regardless of which two nonidentity
members of $D$ are used as its ordered outer generators.
Every such group fixes $e_3$ as a projective point. Its matrix lifts
have bottom eigenvalue $\pm1$, so they meet the scalar centre trivially.
We may work throughout with these actual matrix groups.

Write
\[
 M(x,y)=\begin{pmatrix}1&0&0\\0&1&0\\x&y&1\end{pmatrix}.
\]
A pure block matrix with upper block $B$ and lower scalar
$\delta=\det B\in\{1,-1\}$ acts by conjugation on these rows as
$(x,y)\mapsto\delta(x,y)B$. This elementary action, together with
coordinate sign projections, accounts for the generic cases.

\begin{theorem}\label{generic}
Suppose $bc\ne0$. Normalize
\[
 A=\begin{pmatrix}a&-a-1\\a-1&-a\end{pmatrix},\qquad a\ne\pm1.
                                                               \tag{2}
\]
Then $v$ is either zero or can be normalized to $(1,-1)$.
Put $K=\F_p(a)$, choose $\lambda$ with
$\lambda+\lambda^{-1}=2a$, and set
\[
 h=|\lambda|,\quad h'=|{-\lambda}|,\quad m=\operatorname{lcm}(h,2).
\]
If $v=0$, the generated group is $D_{2m}$. If $v=(1,-1)$, it is
$K^2\rtimes D_{2m}$ in the displayed faithful two-dimensional action.
All three choices of distinct outer generators in $D$ then give full
rank-three IP. Up to reversal their types are
\[
                  [h,h'],\qquad[h,2p],\qquad[h',2p].    \tag{3}
\]
At $a=0$ the group is $D_{2p}\wr C_2$ and these types reduce to
$[4,4]$ and $[4,2p]$.
\end{theorem}
\begin{proof}
Diagonal scaling gives (2), and its left $+1$ eigenspace is
$\gen{(1,-1)}$. Scaling the third coordinate gives the asserted
normalization of $v$. Let $h_A=\operatorname{diag}(A,-1)$.
For $v=(1,-1)$, the actual kernel of the upper-block action contains
\[
                        (t d_3)^2=M(2,-2).
\]
Conjugating by $D$ and adding or subtracting the results separates
the two coordinates. Dividing by two is valid in odd characteristic.
Conjugation by $t$ and projection then multiply a coordinate by $a$.
Iteration gives every element of $K$ in each coordinate. Conversely
all matrices have entries in $K$. Thus this kernel is exactly $M(K^2)$.
Since $t=M(v)h_A$, the group contains the pure block lift $h_A$ and
is the actual semidirect product $M(K^2)\rtimes H$, where
$H=\gen{D,h_A}$. In the zero case it is just $H$.

The upper block of $d_1h_A$ has determinant one and eigenvalues
$\lambda,\lambda^{-1}$. They are distinct and not $\pm1$.
Adjoining $d_3$, whose upper block is $-I$, enlarges its cyclic group
to order $m$: for even $h$, its half-power is already $-I$, whereas
for odd $h$ adjoining $-I$ doubles the order. The remaining reflection
inverts this cyclic group. Hence $H=D_{2m}$. In particular its action
on $K^2$ is faithful and no nonidentity rotation fixes a nonzero vector.
The trace field is $K$; equivalently it is the least field in which
$\lambda+\lambda^{-1}$ lies. Thus $m$ divides $|K|-1$ or $|K|+1$.

For completeness we prove the actual intersections. Put
$A_i=\gen{d_i,t}$, and let $H_i=\gen{d_i,h_A}$ for $i=1,2$.
With the convention $H^g=g^{-1}Hg$, direct multiplication gives
\[
 A_i=H_i^{M(c_i)},\qquad
 c_1=(0,(1-a)^{-1}),\quad c_2=(-(a+1)^{-1},0).          \tag{4}
\]
The translations $c_i$ fix $d_i$ by conjugation and send $h_A$ to $t$.
The vectors differ. Therefore the common elements of these two
complements inject into the stabilizer of $c_1-c_2$ in $H$.
Such a stabilizer has order at most two: a nonidentity rotation fixes
no nonzero vector, and two reflections would have such a rotation
as their product. Since both ends contain $t$, their intersection
is exactly $\gen{t}$.

The third end is
\[
 A_3=M(\F_p v)\rtimes\gen{d_3,h_A},\qquad |A_3|=4p.   \tag{5}
\]
Indeed its squared rotation gives $M(2v)$, which gives $M(v)$ and
hence $h_A$. Its upper-block image is the Klein group
$\gen{d_3,h_A}$. Each $A_i$, $i=1,2$, injects into that block action.
If an element of $A_i\cap A_3$ had image $d_3$, equation (4) would
force its translation to be $-2c_i$. This is not in the line $Kv$,
and hence not in the kernel line of (5). The image $d_3h_A$ is also
excluded by multiplication by the common element $t$. Only the
identity and $h_A$ images remain, so $A_i\cap A_3=\gen{t}$.
The rank-three end criterion now proves full IP for every outer pair.

The rotations $d_1t,d_2t$ have orders $h,h'$: their upper blocks
have no eigenvalue one, so the affine displacement contributes nothing
to their order. Equation (5) gives $|d_3t|=2p$. This proves (3).
If $a=0$, the block group is the signed permutation group on two
coordinates. Its sign changes, together with $\F_p^2$, form two
independent $D_{2p}$ factors, exchanged by its coordinate swap.
This is exactly the asserted wreath product.
\end{proof}

A dihedral group admits a rank-three string C-marking only in the
degenerate case $D_{4k}$ with $k$ odd, of type $[k,2]$ or its reverse.
Here is the short argument needed for the zero case above. If all three
marks are reflections, the commuting ends differ by the central
involution; their common nontrivial squared rotation violates IP unless
an adjacent product has order two. In that remaining situation a middle
reflection commuting with an end is either that end or the other end,
again violating independence. Thus the central involution is an end
mark. The other two marks generate a dihedral subgroup intersecting
it trivially, and the whole group is their direct product with $C_2$.
This is dihedral exactly when their rotation order $k$ is odd.
The middle mark cannot be central, since then the whole group would be
an elementary abelian group of rank at most two.

\section*{The triangular specializations}
If $b=c=0$ in (1), either $t\in D$, which fails IP, or the group is
$D_{4p}$. If exactly one of $b,c$ vanishes, interchanging the first two
coordinates makes $a=1$. Scaling the remaining entry and, when
nonzero, the bottom row gives one of the following matrices, or the
same $D_{4p}$ case:
\[
 t_P=\begin{pmatrix}1&0&0\\1&-1&0\\1&0&-1\end{pmatrix},\qquad
 t_H=\begin{pmatrix}1&-2&0\\0&-1&0\\1&-1&-1\end{pmatrix}.          \tag{6}
\]
For example when $b=0,c\ne0$, the equation $vA=v$ makes $v=(v_1,0)$;
when $c=0,b\ne0$, it makes $v_2=bv_1/2$. This proves completeness
of (6), including the cases with zero bottom row.

\begin{theorem}\label{triangular}
The group $\gen{D,t_P}$ is $D_{2p}\times D_{2p}$. Its two valid
nondegenerate outer pairs have type $[p,2p]$, up to reversal; the
remaining pair fails IP. The group $\gen{D,t_H}$ is
$\operatorname{Heis}(p)\rtimes C_2^2$. All three outer pairs have IP,
with types $[p,2p]$, $[p,2p]$, $[2p,2p]$, up to reversal.
\end{theorem}
\begin{proof}
In the product case write a translation as $n(x,y)=I+xE_{21}+yE_{31}$.
Then $t_P=n(1,1)d_1$. The diagonal signs separate its two prime-field
coordinates: $d_2$ inverts the first, $d_3$ the second, and $d_1$
both. Thus the actual group is the direct product claimed.
The end $A_1=\gen{d_1,t_P}$ has normal forms $n(k,k)d_1^\epsilon$.
In $A_2=\gen{d_2,t_P}$ the coordinate $x$ is arbitrary and $y$ is
zero or one according as its diagonal component fixes or negates the
second coordinate. For $A_3$ the analogous condition holds on $x$.
These descriptions follow by squaring $d_2t_P,d_3t_P$, and give
$A_1\cap A_2=A_1\cap A_3=\gen{t_P}$, whereas $|A_2\cap A_3|=4$.
The orders of $d_it_P$ are respectively $p,2p,2p$.

For the second case use
\[
 U(x,y,z)=I+xE_{12}+yE_{31}+zE_{32},\qquad
 U(x,y,z)U(X,Y,Z)=U(x+X,y+Y,z+Z+yX).
\]
We have $t_H=U(2,1,1)d_1$. Conjugation by $d_2$ changes this
unipotent factor to $U(-2,1,-1)$. Their commutator has central
coordinate $-4$, so it generates the full prime-field central root
group. Their product then supplies the second root group, and the
first follows by division. This proves the actual full Heisenberg
kernel and its diagonal complement $D$.

Put $A_i=\gen{d_i,t_H}$. The first end has forms
$U(2k,k,k^2)d_1^\epsilon$. The kernels of the diagonal projections
of the other two ends are respectively
\[
 N_2=\{U(x,0,x/2):x\in\F_p\},\qquad
 N_3=\{U(0,y,-y):y\in\F_p\}.
\]
These follow by squaring $d_2t_H,d_3t_H$; each end has order $4p$.
The first kernel meets either of them trivially, giving
$A_1\cap A_2=A_1\cap A_3=\gen{t_H}$.
For the remaining intersection the kernels are disjoint. At diagonal
component $d_2$, the possible translations in $A_2,A_3$ are respectively
$U(x,0,x/2)$ and $U(2,t,t)$; equality would force $x=2,t=0$ and
then $1=0$. At component $d_3$ they are $U(x-2,1,x/2-1)$ and
$U(0,y,-y)$; equality would give $x=2,y=1$ and then $0=-1$.
Thus only the identity and $d_1$ components remain, with unique common
lifts $1,t_H$. This proves the final intersection. The rotation orders
are $p,2p,2p$, and the end criterion completes the proof.
\end{proof}

\section*{The irreducible conic alternative}
\begin{theorem}\label{conic}
If an odd-characteristic rank-three string C-group in $\PSL_3(F)$
fixes no projective point or line, it preserves a nonsingular conic.
Its group is one of $S_4,A_5,\PSL_2(K),\PGL_2(K)$, in the natural
three-dimensional orthogonal representation over a subfield $K$.
\end{theorem}
\begin{proof}
With the diagonal outer Klein group, write the middle homology as
$t=-I+2uf$, where $f(u)=1$. Every coordinate of $u$ and $f$ is
nonzero: a zero coordinate of $u$ would preserve the corresponding
coordinate plane, and a zero coordinate of $f$ the coordinate point.
The diagonal quadratic form
\[
                         Q=\operatorname{diag}(f_1/u_1,f_2/u_2,f_3/u_3)
\]
is nonsingular, has $Q(u,u)=1$, and is preserved by $D$ and $t$.
Its determinant-one orthogonal group meets the scalar centre trivially,
and is isomorphic to $\PGL_2(F)$. Apply the retained Dickson theorem.
Semisimple cyclic and dihedral subgroups preserve the torus axis in
this orthogonal representation. The remaining cyclic or dihedral cases
lie in a point stabilizer in the projective-line action, which preserves
its corresponding isotropic point. The same holds for the Borel-type
subgroups in Dickson's list. These are excluded by
hypothesis. An $A_4$ cannot be generated by involutions, whose generated
subgroup is only $V_4$. The stated alternatives remain.
\end{proof}
For $S_4$ the nondegenerate types are $[3,3]$ and $[3,4]$, up to
reversal. Both adjacent orders cannot be four: the involution types in
$S_4$ would give only a subgroup $D_8$. For $A_5$ the possible types
are $[3,5]$ and $[5,5]$: a $[3,3]$ marking would be a quotient of
the Coxeter group $S_4$. The usual tetrahedral and projective cubical
and icosahedral markings supply these small possibilities. The wider
rank-three marking classifications within the projective-line groups
remain separately attributed inputs; we do not repeat them here.

\section*{The source table and the corrected affine row}
Theorems~\ref{generic}--\ref{conic} recover the group alternatives of
Theorems~5.1--5.2. They also give every nondegenerate affine type in
Table~3, with one necessary correction. The generic row in that table
prints only $[k,4]$ and $[k,k]$. Its literal middle matrix (2) with
bottom row $(1,-1,-1)$ instead gives the three types (3).
For a concrete instance take $F=\F_5$ and $a=2$. Then $h=3,h'=6$,
$K=\F_5$, $m=6$, and the group has order $300$. The three actual
C-markings have types $[3,6]$, $[3,10]$, $[6,10]$. Their full subset
intersections are checked in the supplied finite control; the uniform
proof above gives them independently of that computation. This is a
correction to a type row, not a refutation of the subgroup theorem or
of the paper's principal chiral constructions and rank bound.

The centre--axis normalization and original matrices are due to
Leemans--Vandenschrick. The proof here replaces the separate geometric
case calculations by one block action, its translation kernel and
actual complement intersections. No numerical compression or new
conjecture resolution is asserted by this component.
\end{document}
